\slajdid{09_2-s050}
\begin{frame}[label=09_2-s050]{Latch-Up}
\gusttekst\setlength{\abovedisplayskip}{3pt}\setlength{\belowdisplayskip}{3pt}\setlength{\abovedisplayshortskip}{2pt}\setlength{\belowdisplayshortskip}{2pt}Jedan od velikih problema koji se pojavljuje kod CMOS logičkih kola zbog njihove interne strukture i parazitnih efekata jeste pojava lečapa (Latch-Up).\par\medskip
\begin{columns}[c,onlytextwidth]\column{.56\textwidth}\centering\begin{tikzpicture}[x=.72cm,y=.65cm,font=\sitneoznake,>=Latex]
\ctikzset{bipoles/length=.7cm}
\draw(-.4,0)--(7.6,0);\draw(-.4,-3.5)--(7.6,-3.5);\node at(6.9,-3.25){p};
\draw(3.55,0)..controls(3.55,-2.3)and(3.85,-2.3)..(4.15,-2.3)--(6.85,-2.3)..controls(7.4,-2.3)and(7.4,-.6)..(7.4,0);\node at(7,-1.2){n};
\foreach \x/\dop/\lab in{.35/p+/Bn,1.35/n+/Sn,2.9/n+/Dn,4.2/p+/Dp,5.8/p+/Sp,6.85/n+/Bp}{
\draw(\x-.3,0)--(\x-.3,-.17)arc[start angle=180,end angle=360,x radius=.3,y radius=.35]--(\x+.3,0);\node at(\x,-.3){\dop};\draw(\x-.1,0)rectangle(\x+.1,.17);\draw(\x,.17)--(\x,.47)node[above]{\lab};}
\draw(1.75,.03)rectangle(2.45,.17);\draw(2.1,.17)--(2.1,.47)node[above]{Gn};
\draw(4.5,.03)rectangle(5.35,.17);\draw(4.925,.17)--(4.925,.47)node[above]{Gp};
\draw(.35,.84)--(.35,.95)--(-.3,.95)node[ground,rotate=-90]{};\draw(1.35,.84)--(1.35,1.7)--(-.3,1.7)node[ground,rotate=-90]{};
\draw(2.1,.84)--(2.1,2.15)node[ocirc]{}node[above]{$V_I$};\draw(2.9,.84)--(2.9,1.2)--(4.2,1.2)--(4.2,.84);\draw(3.5,1.2)--(3.5,1.45)node[ocirc]{}node[right]{$V_O$};
\draw(4.925,.84)--(4.925,1.9)--(2.1,1.9);\draw(5.8,.84)--(5.8,1.5)--(6.85,1.5)--(6.85,.84);\draw(5.5,1.5)--(7.15,1.5)node[above left]{$V_{DD}$};

\draw(1.65,-2.3)--(2.35,-2.3);\draw(1.8,-2.3)--(1.35,-1.9)--(1.35,-.2);\draw[-{Triangle[length=1.5mm,width=1mm]}](1.72,-2.23)--(1.37,-1.92);\draw(2.2,-2.3)--(2.65,-1.9)--(3.25,-1.9);\draw(2,-2.3)--(2,-2.95);\node[above]at(2,-2.3){$T_N$};
\draw(5.15,-1.55)--(5.15,-2.25);\draw(5.15,-1.7)--(5.8,-1.25)--(5.8,-.2);\draw[-{Triangle[length=1.5mm,width=1mm]}](5.48,-1.47)--(5.17,-1.69);\draw(5.15,-2.1)--(5.8,-2.5)--(5.8,-2.95);\node[above]at(5.05,-1.35){$T_P$};
\draw(.35,-.2)--(.35,-2.95)to[R,l_=$R_{S2}$](2,-2.95)to[R,l_=$R_{S1}$](3.5,-2.95)--(5.8,-2.95);
\draw(3.25,-1.9)to[R](4.75,-1.9)--(5.15,-1.9)to[R,l_=$R_{W1}$](6.85,-1.9)--(6.85,-.2);
\draw(4.2,-.2)--(4.2,-.72);\draw(3.95,-.72)--(4.45,-.72);\draw(3.95,-.92)--(4.45,-.92);\draw(4.2,-.92)--(4.2,-1.35)--(4.75,-1.35)--(4.75,-1.9);\node[left]at(3.95,-.82){$C_P$};
\draw(2.9,-.2)--(2.9,-1.75);\draw[preaction={draw=white,line width=2pt}](2.9,-1.75)arc[start angle=90,end angle=-90,radius=.15];\draw(2.9,-2.05)--(2.9,-2.35);\draw(2.65,-2.35)--(3.15,-2.35);\draw(2.65,-2.55)--(3.15,-2.55);\draw(2.9,-2.55)--(2.9,-2.65)--(2,-2.65);\node[right]at(3.15,-2.65){$C_N$};\node at(4.75,-2.65){$R_{W2}$};
\end{tikzpicture}


\column{.42\textwidth}
Dimenzije tranzistora su male i dovoljne da se formiraju dva parazitna tranzistora u npn i pnp oblastima. Ovi tranzistori imaju „usku“ bazu i mogu da rade sa relativno velikim pojačanjima.
\begin{center}\begin{tikzpicture}[x=.6cm,y=.55cm,font=\sitneoznake,>=Latex]
\ctikzset{bipoles/length=.7cm}
\draw(2.35,6)--(4.3,6);\node[above]at(3.3,6){$V_{DD}$};
\draw(1.8,4.3)--(1.8,5.1);\draw(1.8,4.9)--(2.55,5.25)--(2.55,6);\draw[-{Triangle[length=1.7mm,width=1.2mm]}](2.3,5.13)--(1.82,4.91);\draw(1.8,4.5)--(2.55,4.15)--(2.55,3.45);\node[right]at(2.55,5.3){$T_P$};
\draw(3.3,1.7)--(3.3,2.5);\draw(3.3,2.3)--(4.05,2.65)--(4.05,3.25);\draw(3.3,1.9)--(4.05,1.5)--(4.05,.65)node[ground]{};\draw[-{Triangle[length=1.7mm,width=1.2mm]}](3.56,1.76)--(4.03,1.51);\node[right]at(4.05,2.25){$T_N$};
\draw(4.05,6)to[R,l=$R_{W1}$](4.05,4.7)to[R,l=$R_{W2}$](4.05,3.25);
\draw(2.55,3.45)to[R,l_=$R_{S1}$](2.55,2.1)to[R,l_=$R_{S2}$](2.55,.65)node[ground]{};\draw(2.55,2.1)--(3.3,2.1);
\draw(.35,4.7)--(.8,4.7);\draw(.8,4.4)--(.8,5);\draw(1,4.4)--(1,5);\draw(1,4.7)--(4.05,4.7);\node[above]at(.9,5.04){$C_P$};
\draw(.35,2.1)--(.8,2.1);\draw(.8,1.8)--(.8,2.4);\draw(1,1.8)--(1,2.4);\draw(1,2.1)--(2.55,2.1);\node[left]at(.65,1.65){$C_N$};
\draw(.35,2.1)--(.35,4.7);\draw(-.25,3.35)node[ocirc]{}node[below left]{$V_O$}--(.35,3.35);
\draw(-.95,3.75)--(-.65,3.75)--(-.65,4.05)--(-.3,4.05)--(-.3,3.75)--(.0,3.75);
\begin{scope}[shift={(-.1,-.25)}]\draw(1.17,4.4)--(1.17,4.66)..controls(1.2,4.4)and(1.25,4.4)..(1.43,4.4)--(1.58,4.4)--(1.58,4.11)..controls(1.61,4.4)and(1.67,4.4)..(1.78,4.4);\end{scope}
\begin{scope}[shift={(-.4,0)}]\draw(1.17,2.55)--(1.17,2.81)..controls(1.2,2.55)and(1.25,2.55)..(1.43,2.55)--(1.58,2.55)--(1.58,2.26)..controls(1.61,2.55)and(1.67,2.55)..(1.78,2.55);\end{scope}
\end{tikzpicture}
\end{center}\end{columns}
\end{frame}
